\section{Duality shifts and the cylinder signs}
\label{app:duality-signs}

This appendix fixes one coherent convention and verifies the signs used
in \Cref{cat:pairing-form,cat:endpoint-forms}.  It also explains why an
overall sign in the cylinder does not alter homotopy invariance.  The
translated-duality convention is the one in
\cite[\S2.1]{Woolf2008}.  For exactness sign $+1$, the cone identities
are those of \cite[Lemma~2.1]{Woolf2008}; their version for
$\delta\in\{1,-1\}$ is proved in \Cref{cat:symmetric-cone}.
The calculations below apply to complexes
of sheaves by using resolutions for internal Hom.  They involve only
the signs of complexes and therefore commute with restriction and with
the natural derived-functor comparisons used in the text.

\subsection{Hom, shifts, and biduality}

For a cochain complex $A$, put
\[
 A[j]^i=A^{i+j},\qquad d_{A[j]}=(-1)^j d_A.
\]
If $a$ has degree $i$ in $A$, write $a[j]$ for the corresponding element
of degree $i-j$ in $A[j]$.  For a fixed complex $K$, the Hom complex
$D A=\mathcal Hom(A,K)$ has differential
\[
 d(F)=d_KF-(-1)^nFd_A\qquad(|F|=n).
\]
The shift comparison is
\begin{equation}
 \label{cat:hom-shift}
 \begin{split}
 \kappa_j:D(A[j])&\longrightarrow(DA)[-j],\\
 \kappa_j(F)(a)&=(-1)^{jn-j(j-1)/2}F(a[j]),
 \qquad F\in D(A[j])^n.
 \end{split}
\end{equation}
For $j=1$ this is the factor $(-1)^n$.  To check that
\eqref{cat:hom-shift} is a chain map, note that increasing $n$ by one
changes its exponent by $j$, precisely the differential sign of the
target shift.  The other differential term has the extra $(-1)^j$
from the source shift.  These are the two terms in the Hom
differential, so both commute with $\kappa_j$.  Iterating the
one-step comparison gives the exponent
$n+(n-1)+\cdots+(n-j+1)$ for positive $j$; the formula for negative
$j$ follows by inversion.  We use these iterated comparisons whenever
dual functors and shifts are identified.

The unshifted biduality is the tensor-symmetry evaluation
\begin{equation}
 \label{cat:hom-biduality}
 \eta_A(a)(F)=(-1)^{|a||F|}F(a).
\end{equation}
It is a chain map, as follows by substituting the two Hom differentials;
its two evaluation signs also give
$D(\eta_A)\eta_{DA}=\id_{DA}$ whenever the duality is defined.
The comparisons \eqref{cat:hom-shift} make this unshifted dual a
$+1$-exact duality: dualizing a cone triangle and identifying its shifts
gives the dual triangle with positive final dual arrow.

Now identify $D_j=[j]D$ with $\mathcal Hom(-,K[j])$ by the output-shift
comparison $f[j](a)=f(a)[j]$, with no additional sign. Its two
differential expressions both equal
$(-1)^j d_Kf-(-1)^n f d_A$ for an element of shifted degree $n$.
Tensor symmetry into $K[j]$ gives the shifted biduality
$\eta^{(j)}$.  Under the iterated comparisons of
\eqref{cat:hom-shift}, its expression relative to the unshifted
biduality is
\begin{equation}
 \label{cat:translated-biduality}
 \eta^{(j)}=(-1)^{j(j+1)/2}\eta.
\end{equation}
Indeed, on an element $a$ of degree $i$, moving the inner $j$-shift
through the outer Hom contributes
$(-1)^{ji+j(j+1)/2}$.  In the evaluation
\eqref{cat:hom-biduality}, a dual element now has its degree changed
by $j$, contributing the additional factor $(-1)^{ji}$.  The two
terms depending on $i$ cancel, leaving
\eqref{cat:translated-biduality}.  The exactness sign of $D_j$ is
$(-1)^j$.

It follows directly from tensor--Hom adjunction that a morphism
$A\to\mathcal Hom(A,K[j])$ is symmetric for this biduality precisely
when its pairing $A\otimes A\to K[j]$ is invariant under Koszul
interchange.  Apply this with $K=\omega_X$ and $j=-r$.  This proves
\eqref{cat:pairing-form} and its agreement with $W_r=W^{-r}$.
In particular, if $X$ is an oriented $d$-manifold, the target
$\omega_X[-d]$ is the constant sheaf in degree zero.  The pairing
$1\otimes1\mapsto1$ is symmetric and positive in an oriented local
trivialization.  No additional parity sign is attached to this
orientation form.

\subsection{Evaluation of a product}

Fix the convention for combining output shifts by
\begin{equation}
 \label{cat:tensor-output-shifts}
 K[p]\otimes L[q]\longrightarrow(K\otimes L)[p+q],
 \qquad k[p]\otimes\ell[q]\longmapsto
 (-1)^{q|k|}(k\otimes\ell)[p+q],
\end{equation}
where $|k|$ is the degree in the unshifted $K$.  This is a chain map:
in the first differential term the extra exponent is $q$, and in the
second the two exponents differ by $2p$.  This choice is associative,
since the exponents for three factors in either order are
$q|k|+t|k|+t|\ell|$.

Let $f,g$ have degrees $l,n$ in $\D_{I,s}B,\D_{X,r}A$ respectively.
Thus their degrees before the output shifts are $l-s,n-r$.
For homogeneous $b,a$, the product map defined in
\Cref{cat:interval-products} evaluates as
\begin{equation}
 \label{cat:product-evaluation-sign}
 \mu_{s,r}(f\boxtimes g)(b\boxtimes a)
 =(-1)^{n|b|+r(|b|+l-s)}f(b)\boxtimes g(a),
\end{equation}
followed by the product trace.  The exponent $n|b|$ moves $g$ past
$b$.  The first unshifted output has degree $|b|+l-s$; applying
\eqref{cat:tensor-output-shifts} to output shifts $-s,-r$ gives the
second exponent.  This derives the formula without assuming any
finiteness or local constancy of the stalks of $A$.

\begin{lemma}[Suspending the first product input]
\label{cat:product-shift-sign}
Let $B\in\mathcal C_I$, $A\in\T(X)$ (or the corresponding
semialgebraic category), and $s,r\in\Z$. Write
\[
 \gamma_s(B)=[-s]\kappa_1:
 \D_{I,s}(B[1])\longrightarrow(\D_{I,s}B)[-1].
\]
Use the first-factor shift comparisons
\[
 \begin{aligned}
 \sigma: B[1]\boxtimes A&\xrightarrow{\sim}(B\boxtimes A)[1],\\
 \epsilon: (\D_{I,s}B)[-1]\boxtimes\D_{X,r}A
 &\xrightarrow{\sim}(\D_{I,s}B\boxtimes\D_{X,r}A)[-1].
 \end{aligned}
\]
They introduce no sign under \eqref{cat:tensor-output-shifts}.
Indicate the inputs of the product comparison by superscripts and put
\[
 \widetilde\mu_{s,r}
   =(\D_{I\times X,s+r}\sigma)^{-1}\mu_{s,r}^{B[1],A},
 \qquad
 \lambda=\epsilon\circ(\gamma_s(B)\boxtimes1).
\]
Then the two routes to the suspended total dual satisfy
\begin{equation}\label{cat:suspension-comparison}
 \gamma_{s+r}(B\boxtimes A)\,\widetilde\mu_{s,r}
   =(-1)^r\mu_{s,r}^{B,A}[-1]\,\lambda.
\end{equation}
The scalar $(-1)^r$ is independent of $A$, its cohomological degrees,
and the point of $I$.
\end{lemma}

\begin{proof}
Keep $l,n$ as the degrees of the two shifted dual elements.  The Hom
shift comparison on the first factor contributes $(-1)^{l-s}$, while
the comparison on the total dual contributes
$(-1)^{l+n-s-r}$.  In
\eqref{cat:product-evaluation-sign}, replacing $b$ by $b[1]$ changes
the braiding exponent by $-n$.  After the first Hom comparison, the
degree of the corresponding dual element changes along with that of
$b$, so the output-shift term is the same along the two routes.
The discrepancy is therefore
\[
 (l+n-s-r)-(l-s)-n\equiv r\pmod2.
\]
All dependence on the individual degrees cancels.  This is the stated
scalar, and it is the suspension discrepancy used when dualizing
$\partial\boxtimes1$ in \Cref{cat:endpoint-forms}.
\end{proof}

\subsection{Normalizing the endpoint calculation}

Recall $H=\const{\R}{I}$, $J=j_!\const{\R}{(0,1)}$, and
$E=E_0\oplus E_1$ from \eqref{cat:interval-objects}. The interval
triangle is
\[
 J\xrightarrow{u}H\xrightarrow{v}E\xrightarrow{\partial}J[1].
\]
Choose the
orientation identification so that its connecting homomorphism is
the difference map in \eqref{cat:interval-cohomology}.  Proper duality
to the point identifies its shifted dual with
\[
 [1](\D_{I,1}\partial):H\longrightarrow E,
 \qquad h\longmapsto(-h|_0,h|_1).
\]
The symmetric-cone identity
$\psi_Ev=-[1](\D_{I,1}\partial)$ therefore gives
\begin{equation}
 \label{cat:interval-cone-matrix}
 \psi_E=
 \begin{pmatrix}1&0\\0&-1\end{pmatrix}
 :E\longrightarrow\D_I E.
\end{equation}
There are no off-diagonal maps because the endpoint supports are
disjoint.  The diagonal coefficients are determined by precomposition
with $v:H\to E$, since on each endpoint this restriction is the
identity.  Thus the cohomology calculation determines the full
derived-category morphism.

Now let $(A,\alpha)$ be a degree-$r$ form on $X$, and write
$i_t:X\to I\times X$ for the endpoint inclusion, $t=0,1$.
For the product cone of \eqref{cat:cylinder-morphism}, the same cone
identity, the naturality of \eqref{cat:external-duality}, and
\Cref{cat:product-shift-sign} give
\begin{equation}
 \label{cat:cylinder-cone-matrix}
 \psi_r=(-1)^r
 \begin{pmatrix}(i_0)_*\alpha&0\\0&-(i_1)_*\alpha\end{pmatrix}.
\end{equation}
One can check this equality by precomposing with
$v\boxtimes1$, which determines each diagonal as in
\Cref{cat:endpoint-forms}.  Combining the shifted interval and total
dualities uses \Cref{cat:product-shift-sign}; desuspending the first
dual output thereafter leaves \eqref{cat:tensor-output-shifts}
unchanged.  There is consequently just the displayed factor
$(-1)^r$.

Reversing the orientation or changing all cone normalizations
coherently may negate both entries of
\eqref{cat:cylinder-cone-matrix}.  It leaves their relative sign
opposite.  Since the form is neutral, both normalizations give the
same equality of the two endpoint Witt classes.  This is the only
sign consequence used in \Cref{thm:homotopy}.
